% fig_inverse_falsifier.tex -- native pgfplots (figure*, full width).
% The negative-order (inverse) f-sum moment m_{-1}=int A_J(omega)/omega domega as a LOW-frequency falsifier.
% Flagship doped L=12 Hubbard ring (PBC, 2/3 filling, U/t=8), sector-Lanczos current spectral function.
% Data: run_inverse_moment_falsifier.py -> inverse_falsifier_spectrum.dat, inverse_falsifier_sweep.dat.
\begin{tikzpicture}
\begin{groupplot}[
  group style={group size=2 by 1, horizontal sep=2.7cm},
  paperaxis, width=8.3cm, height=6.2cm,
]
% ---- panel (a): the spectrum and the 1/omega falsifier kernel ----
\nextgroupplot[
  xlabel={$\omega$ \; (energy above ground state, $t$)},
  ylabel={current spectral weight $A_J(\omega)$},
  xmin=0, xmax=15, ymin=0, ymax=1.42,
  ytick={0,0.5,1},
  title={\small\textbf{(a)} the inverse kernel $1/\omega$ weights the depleted low-$\omega$ region},
  clip=false,
]
  % the 1/omega falsifier kernel (rescaled to the axis), shaded: shows where m_{-1} is most sensitive
  \addplot[pTeal, densely dotted, line width=0.9pt, fill=pTeal, fill opacity=0.06]
      table[x=omega,y=kernel]{figs/inverse_falsifier_spectrum.dat} \closedcycle;
  % true spectrum: the mid-IR / Mott band carries essentially all the regular weight (low-omega empty)
  \addplot[draw=pIndigo, thick, fill=pIndigo, fill opacity=0.15]
      table[x=omega,y=Atrue]{figs/inverse_falsifier_spectrum.dat} \closedcycle;
  % wrong reconstruction: a spurious low-omega (Drude-like) peak that steals mid-IR weight (equal m_0)
  \addplot[pOx, very thick] table[x=omega,y=Awrong]{figs/inverse_falsifier_spectrum.dat};
  \node[pIndigo, font=\small, anchor=north east, align=right] at (axis cs:14.6,1.41)
      {true $A_J(\omega)$\\ (mid-IR band)};
  \node[pOx, font=\small, anchor=west, align=left] at (axis cs:3.35,1.20) {spurious low-$\omega$\\ (Drude-like) peak};
  \draw[pOx, -{Stealth[length=1.4mm]}, line width=0.5pt] (axis cs:3.20,1.12) -- (axis cs:2.75,1.20);
  \node[pTeal, font=\scriptsize, anchor=west, align=left] at (axis cs:4.2,0.66) {falsifier kernel\\ $1/\omega$};
  \draw[pTeal, -{Stealth[length=1.2mm]}, line width=0.4pt] (axis cs:4.3,0.50) -- (axis cs:4.3,0.13);
% ---- panel (b): relative discrepancies vs misplaced weight fraction ----
\nextgroupplot[
  xlabel={weight fraction misplaced to spurious low-$\omega$ peak (\%)},
  ylabel={relative moment discrepancy $\Delta_k/m_k$ (\%)},
  xmin=0, xmax=50, ymin=-20, ymax=165,
  ytick={0,50,100,150},
  title={\small\textbf{(b)} the inverse moment moves ${\sim}4\times$ more than $m_1$},
  clip=false,
]
  \addplot[threshold, forget plot] coordinates {(0,5) (50,5)};
  % inverse moment: fires hardest
  \addplot[firesline] table[x=f,y=dm1]{figs/inverse_falsifier_sweep.dat};
  % positive moments: modest
  \addplot[pIndigo, very thick, mark=triangle*, mark options={fill=pIndigo}]
      table[x=f,y=d2]{figs/inverse_falsifier_sweep.dat};
  \addplot[pMuted, thick, mark=diamond*, mark options={fill=white, draw=pMuted}]
      table[x=f,y=d1]{figs/inverse_falsifier_sweep.dat};
  % zeroth moment: blind (flat on zero)
  \addplot[blindline] table[x=f,y=d0]{figs/inverse_falsifier_sweep.dat};
  % inverse-moment label: upper-left empty wedge, above the steep inverse line
  \node[pOx, font=\small, anchor=south west, align=left] at (axis cs:2,132)
      {inverse moment $m_{-1}$\\ ($\Delta_{-1}/m_{-1}$)};
  % positive-moment labels: at the right ends of their (shallow) lines
  \node[pIndigo, font=\scriptsize, anchor=south east] at (axis cs:44,49) {$m_2$};
  \node[pMuted, font=\scriptsize, anchor=north west] at (axis cs:44,30) {$m_1$};
  % threshold label: right gap between the y=5 line and the shallow m_1 line
  \node[pMuted, font=\scriptsize, anchor=west] at (axis cs:24,14) {$5\%$ threshold};
  % blind label: centered in the widened empty strip below the m_0 line, clear of the x-axis
  \node[pBlind, font=\small] at (axis cs:25,-11)
      {total weight $m_0=\langle J^2\rangle$ (blind)};
\end{groupplot}
\end{tikzpicture}
